An integer $x$ is impossible since then $x^2+b>0$. Also $x^2=a\{x\}-b<9$, so $m=\lfloor x\rfloor\in\{-3,-2,-1,0,1,2\}$. Put $t=x-m\in(0,1)$. The equation becomes $F(t)=t^2+(2m-a)t+m^2+b=0$, with $F(0)>0$. A root in $(0,1)$ occurs exactly when $F(1)<0$, that is $a-b>(m+1)^2$. To justify this equivalence, if the vertex of this monic quadratic lies strictly between zero and one, its integer linear coefficient must be $-1$, and its minimum is $m^2+b-\dfrac{1}{4}>0$. Otherwise it is monotone on this interval or has its minimum outside, so crossing requires opposite endpoint signs. Equality at the right endpoint only gives the excluded integer. For the stated inequality there is exactly one interior root. Therefore all solutions are
\[\boxed{x=\dfrac{a-\sqrt{a^2-4am-4b}}{2},\quad 1\le a,b\le9,\quad m\in\{-3,-2,-1,0,1\},\quad a-b>(m+1)^2.}\]
The strict digit inequalities constitute an explicit finite parametrization, with no additional root filter required.
